% Source: https://artofproblemsolving.com/community/c3h1362069_perfectabundantdeficient_numbers
% Stable identity: aops-topic-1362069-q01
If $n=p_1^{a_1}\cdots p_k^{a_k}$, then
\[
\sigma(n)=\prod_{i=1}^{k}\left(1+p_i+\cdots+p_i^{a_i}\right)=\prod_{i=1}^{k}\dfrac{p_i^{a_i+1}-1}{p_i-1},
\]
because expanding the product produces every divisor exactly once. It is often quicker to compare $\dfrac{\sigma(n)}{n}=\prod\dfrac{\sigma(p_i^{a_i})}{p_i^{a_i}}$ with $2$.

$496=2^4\cdot31$: $\sigma=31\cdot32=992=2\cdot496$, so $496$ is \textbf{perfect}.

$945=3^3\cdot5\cdot7$: $\sigma=40\cdot6\cdot8=1920>1890=2\cdot945$, so $945$ is \textbf{abundant} (it is the smallest odd abundant number).

$1000=2^3\cdot5^3$: $\sigma=15\cdot156=2340>2000$, so $1000$ is \textbf{abundant}.

So: $\boxed{496\text{ perfect},\ 945\text{ abundant},\ 1000\text{ abundant}}$.
